\section{Prescribed spanning paths and arbitrary endpoints}
\label{sec:paths}
A Hamilton cycle does not itself give a Hamilton path with
arbitrarily specified endpoints.
We verify the stronger spanning-path input separately.

We use Theorem 6.4 and Definition 3.9 of~\cite{BJ}.
For a regular balanced bipartite graph of degree \(d\), their
bi-uniformity condition for \(S\) in one part is
\[
 \left||N(v)\cap S|-\frac{2|S|d}{N}\right|\leq\alpha d
 \quad\hbox{at every vertex of the other part}.              \tag{41}
\]
The theorem supplies internally disjoint paths joining prescribed
opposite-part pairs, with all remaining vertices in their interiors,
provided the endpoint sets and multisets obey (41), have equal
cardinality, and the following conditions hold:
\[
 0<c<1/4,\quad 0<\gamma<(\log N)^{-2},\quad
 d'\leq d^{1-c},\quad\alpha\leq d^{-c},\quad
 m\leq d^{-c/900}N,\quad d\geq\gamma^{-10^6/c}.               \tag{42}
\]
Here the graph is \((d\pm d')\)-nearly regular and a
\(\gamma\)-expander in the sense used in (2), and \(N\)
is sufficiently large depending on \(c\).
We apply it to exact regular factors, so \(d'=0\).

\begin{lemma}[Arbitrary small endpoint sets meet the source condition]
\label{lem:smallendpoints}
Let \(H\) be a simple \(d\)-regular balanced bipartite graph
on \(N\) vertices.
Any subset \(S\) of one part of cardinality \(m\) obeys
(41) with \(\alpha=m/d\).
The same holds when its distinct elements are viewed as a multiset.
\end{lemma}
\begin{proof}
Both \(|N(v)\cap S|\) and \(2md/N\) lie in \([0,m]\),
since \(d\leq N/2\). Their difference has modulus at most \(m\).
This is (41), including the set and multiset versions.
\end{proof}

\begin{theorem}[Almost-linear prescribed spanning linkages]
\label{thm:linkages}
Fix \(q,R\) as in Theorem~\ref{thm:filterclassification},
with \(g_R=2\), and fix \(0<\eta<1\).
For all sufficiently large even \(N\), choose any integer
\(1\leq m\leq N^{1-\eta}\), any distinct odd vertices
\(a_1,\ldots,a_m\), and any distinct even vertices
\(b_1,\ldots,b_m\).
Then \(G_N(q,R)\) contains pairwise vertex-disjoint paths
\(P_1,\ldots,P_m\), where \(P_i\) joins \(a_i\) to \(b_i\)
and their vertex sets partition all of \([1,N]\).
In particular it has a Hamilton path between every prescribed
odd-even pair.
\end{theorem}
\begin{proof}
Use the exact regular factor from
Theorem~\ref{thm:filterclassification}.
Its degree is \(d=(1-o(1))\sigma_R N/\log N\),
and its fixed positive spectral gap supplies a fixed positive
edge-expansion constant \(\gamma_0\), under the definition (2).
Choose
\[
 c=\min\{1/5,\eta/2\},\qquad
 \gamma=(\log N)^{-3},\qquad d'=0,\qquad\alpha=m/d .
\]
For large \(N\), \(\gamma<\gamma_0\), so expansion with this
smaller constant holds throughout the required cut range.
The arbitrary endpoint sets meet (41), by
Lemma~\ref{lem:smallendpoints}, and each has cardinality \(m\).
Their vertices are distinct, so the multiset requirements are
the same.

We now check every quantitative condition in (42).
The choice of \(c\) is in \((0,1/4)\), and the choice of
\(\gamma\) is smaller than \((\log N)^{-2}\).
Also
\[
 m/d^{1-c}
 =O_{q,R}\left(N^{c-\eta}(\log N)^{1-c}\right)\longrightarrow0,
\]
since \(c<\eta\).
Thus \(\alpha=m/d\leq d^{-c}\).
Similarly
\[
 \frac{m}{d^{-c/900}N}
 =O_{q,R}\left(N^{-\eta+c/900}
                    (\log N)^{-c/900}\right)\longrightarrow0.
\]
Finally
\[
 d\asymp_{q,R}N/\log N
       \gg(\log N)^{3\cdot10^6/c}
       =\gamma^{-10^6/c}.
\]
The degree error \(d'=0\) obeys its condition exactly.
The cited theorem gives the desired paths.
Their interiors avoid every prescribed endpoint; since all
endpoints are distinct, internal disjointness becomes full
vertex disjointness, and its covering conclusion partitions
the entire vertex set.
\end{proof}

\begin{corollary}[All parity-compatible prime-sum endpoints]
\label{cor:oddendpoints}
For all sufficiently large \(N\), in the ordinary prime-sum
graph on \([1,N]\), every distinct pair \(a,b\) satisfying
the necessary parity condition is joined by a Hamilton path:
opposite parities for even \(N\), and both odd for odd \(N\).
\end{corollary}
\begin{proof}
The even case is Theorem~\ref{thm:linkages} with \(q=2\)
and \(R=\{1\}\).
Suppose \(N\) is odd and \(a,b\) are distinct odd vertices.
Choose an even neighbour \(v\) of \(a\) in \(G_N\).
Such a neighbour exists for every large \(N\), uniformly in
\(a\): the degree estimate (26) for \(G_{N+1}\) loses at
most one edge on deleting vertex \(N+1\).

Consider \(G_N-a\), a balanced bipartite graph on \(N-1\)
vertices. It is obtained from \(G_{N+1}\) by deleting the
odd vertex \(a\) and the even vertex \(N+1\).
Its degrees remain \((1+o(1))(N+1)/\log(N+1)\),
uniformly in the two deleted vertices.
If \(W\) is a subset of the remaining vertices of size at
most \((N-1)/2\), the cut bound (28) in \(G_{N+1}\) loses
at most \(2|W|\) edges. Thus its boundary is at least
\[
 \left(\frac15\frac{N+1}{\log(N+1)}-2\right)|W|.
\]
The complementary smaller-side bound is identical.
The regular-factor lemma supplies a regular spanning factor
of degree \(d=(1-o(1))(N+1)/\log(N+1)\), discarding only
\(o(d)\) edges per vertex.
It retains a fixed positive edge-expansion constant.
The argument proving Theorem~\ref{thm:linkages}, or (42)
with \(m=1\), supplies a Hamilton path from \(v\) to \(b\)
in this factor. Prepend the edge \(av\).
The resulting path begins at \(a\), ends at \(b\), and
covers exactly all \(N\) vertices of \(G_N\).
\end{proof}

\begin{corollary}[Diffuse kernel supports admit prescribed spanning linkages]
\label{cor:weightedlinkages}
Fix \(0<a\leq1\), \(C>0\), \(0<\delta<1\), and
\(0<\zeta<a\).
For all sufficiently large \(N\), every stochastic bipartite
kernel \(P\) with nontrivial gap at least \(\delta\) and
\(\alpha\leq C N^{-a}\) has the following property:
any \(m\leq N^{a-\zeta}\) prescribed distinct opposite-part
pairs can be joined by vertex-disjoint paths partitioning
all \(N\) vertices of its support.
\end{corollary}
\begin{proof}
In Theorem~\ref{thm:rounding}, choose
\[
 d=\lfloor N^{a-\zeta/2}\rfloor.
\]
For all large \(N\), \(d\leq N/2\) and
\(d\alpha L^2=O(N^{-\zeta/2}L^2)\to0\).
The entire error in (1) tends to zero, so its simple regular
factor has nontrivial gap at least \(\delta/2\).
Choose
\[
 c=\min\{1/5,\zeta/(4a)\},\qquad\gamma=(\log N)^{-3}.
\]
As before take \(d'=0\), \(\alpha'=m/d\).
The arbitrary endpoint sets satisfy (41) with \(\alpha'\).
Since
\[
 (a-\zeta/2)(1-c)>(a-\zeta),
\]
we have \(m\leq d^{1-c}\) for all large \(N\).
Moreover
\[
 a-\zeta+(a-\zeta/2)c/900
 \leq 1-\zeta+\zeta/3600<1,
\]
so \(m\leq d^{-c/900}N\).
The polynomial degree dominates
\(\gamma^{-10^6/c}\).
Every hypothesis of (42) is satisfied.
The same spanning-path theorem therefore applies.
\end{proof}
